\subsection{CT6 --- Active/dormant dichotomy}\label{ct:6}\label{the-activedormant-dichotomy}
\ctsignature{\(P_{1\to6}\to \ctcert{C1};\ \ctint{P_{6\to4},P_{6\to3},P_{6\to7},P_{6\to9}};\ \ctrec{P_{6\to10}}\).}

CT6 is an exact sequential machine for the active/dormant split.  Its stable
entry interface fixes the ambient object, branch state, monitored family,
activity threshold, first-failure order, first-failure witness type, scope
predicate, and C1 claim.  A validated \texttt{Input} contains the monitor,
threshold, and order.  The executable \texttt{Framework} adds the predicates
certified by the machine and the exact entry contracts of CT3, CT4, CT7, CT9,
and CT10.

The entry contract is kept in \texttt{CT6.Interface}, independently of the
machine implementation.  This separation is necessary because the tactic
graph contains cycles, while Lean imports must be acyclic.  A producer may
therefore construct an exact \texttt{CT6.Interface.Input} without importing
CT6's nodes or any tactic to which CT6 routes.  The same interface layer is
used for CT7--CT10.

The core plan has five independent fields: \texttt{scope},
\texttt{definition}, \texttt{activity}, \texttt{activeLedger}, and
\texttt{dormant}.  The scope node returns a typed \texttt{ScopeCandidate} or
\texttt{ScopedState}.  The definition node has one successor and certifies
the activity predicate, threshold, order, and first-failure extractor.  The
activity node returns either an indexed \texttt{ActiveState} or an indexed
\texttt{DormantState}.  A dormant state contains the actual first-failure
witness and its structural proof.

The active branch has no further classification choice: the
\texttt{activeLedger} certification node constructs the exact CT4 input and
its alignment proof.  The dormant node returns exactly one of C1, CT3, CT7,
CT9, or CT10.  Each payload is indexed by the dormant state that licensed it
and contains the consumer's validated input.  \texttt{Port} and
\texttt{HandoffPlan} certify consumer acceptance after \texttt{runCore}; they
cannot change the route selected by a core node.

The soundness contracts occur at the following machine boundaries.
\begin{itemize}
\tightlist
\item \textbf{S-Def} fixes the monitor and first-failure vocabulary in the entry interface and is certified before the activity decision.
\item \textbf{S-Dich} is the exhaustive active/dormant decision and the exhaustive dormant routing decision.
\item \textbf{S-Equiv} is carried by the structural first-failure certificate used on the dormant branch.
\item \textbf{S-Rout} and \textbf{S-Trig} are carried by the exact CT3, CT4, CT7, CT9, and CT10 inputs and their alignment proofs.
\end{itemize}

Failure to define activity or to prove the first-failure extractor structural
is a pre-admission defect: no \texttt{DefinitionState}, and hence no admitted
\texttt{CorePlan}, exists.  Absence of any bounded first-failure interface is
the typed scope terminal.  It is not conflated with the CT10 recovery, which
requires a concrete bounded witness and an exact refinement input.

\begin{figure}[p]
\centering
\vspace*{-1.70cm}
\resizebox{\textwidth}{!}{%
\begin{tikzpicture}[x=1cm,y=1cm,every node/.style={font=\scriptsize}]
\node[bcwide] (entry) at (0,0) {{\tiny\ttfamily CT6.entry}\\\textbf{Validated monitor}\\threshold and first-failure order};
\node[bcdec,bclemma,text width=3.45cm,minimum width=4.20cm] (scope) at (0,-2.25) {{\tiny\ttfamily CT6.decide.scope}\\\textbf{Bounded first-failure interface?}};
\node[bcscopefail,text width=3.15cm,minimum width=3.15cm] (scopeout) at (7.10,-2.25) {{\tiny\ttfamily CT6.terminal.scope}\\\textbf{Scope exit}};
\node[bcwide,bclemma] (definition) at (0,-4.55) {{\tiny\ttfamily CT6.certify.definition}\\\textbf{S-Def/S-Equiv certificate}\\activity and first-failure data};
\node[bcdec,bclemma,text width=3.45cm,minimum width=4.20cm] (activity) at (0,-6.95) {{\tiny\ttfamily CT6.decide.activity}\\\textbf{Active or dormant?}};
\node[bcwide,bclemma,text width=3.55cm,minimum width=3.55cm] (ledger) at (7.10,-6.95) {{\tiny\ttfamily CT6.certify.activeLedger}\\\textbf{Construct active ledger}};
\node[bcroute,bcrouteneutral,bcoutputroute,text width=3.05cm,minimum width=3.05cm] (ct4) at (7.10,-9.55) {{\tiny\ttfamily CT6.terminal.ct4}\\\textbf{CT 4 exit}};
\node[bcroutechoice,bctheorem,text width=3.30cm,minimum width=4.20cm] (dormant) at (0,-9.75) {{\tiny\ttfamily CT6.decide.dormant}\\\textbf{Classify structural first failure}};
\node[bcclosed,text width=2.80cm,minimum width=2.80cm] (c1) at (-9.20,-13.00) {{\tiny\ttfamily CT6.terminal.c1}\\\textbf{C1}};
\node[bcroute,bcrouteneutral,bcoutputroute,text width=2.85cm,minimum width=2.85cm] (ct3) at (-4.65,-13.00) {{\tiny\ttfamily CT6.terminal.ct3}\\\textbf{CT 3 exit}};
\node[bcroute,bcrouteneutral,bcoutputroute,text width=2.85cm,minimum width=2.85cm] (ct7) at (0,-13.00) {{\tiny\ttfamily CT6.terminal.ct7}\\\textbf{CT 7 exit}};
\node[bcroute,bcrouteneutral,bcoutputroute,text width=2.85cm,minimum width=2.85cm] (ct9) at (4.65,-13.00) {{\tiny\ttfamily CT6.terminal.ct9}\\\textbf{CT 9 exit}};
\node[bcroute,bcrouteneutral,bcdesign,text width=2.95cm,minimum width=2.95cm] (ct10) at (9.30,-13.00) {{\tiny\ttfamily CT6.terminal.ct10}\\\textbf{CT 10 exit}};
\draw[bcarr] (entry)--(scope);
\draw[bcscopearr] (scope)--node[bclabel,above]{exit}(scopeout);
\draw[bcarr] (scope)--node[bclabel,right]{ready}(definition);
\draw[bcarr] (definition)--node[bclabel,right]{certified}(activity);
\draw[bcarr] (activity)--node[bclabel,above]{active}(ledger);
\draw[bchandoff] (ledger)--node[bclabel,right]{\(P_{6\to4}\)}(ct4);
\draw[bcarr] (activity)--node[bclabel,right]{dormant}(dormant);
\draw[bcarr] (dormant)--node[bclabel,above,sloped]{target hit}(c1);
\draw[bchandoff] (dormant)--node[bclabel,above,sloped]{\(P_{6\to3}\)}(ct3);
\draw[bchandoff] (dormant)--node[bclabel,right]{\(P_{6\to7}\)}(ct7);
\draw[bchandoff] (dormant)--node[bclabel,above,sloped]{\(P_{6\to9}\)}(ct9);
\draw[bcrecoverarr] (dormant)--node[bclabel,above,sloped]{\(P_{6\to10}\)}(ct10);
\end{tikzpicture}%
}
\caption[Closure-tactic 6: exact active/dormant machine]{Closure-tactic 6 as an exact sequential machine.  Monospaced labels are shared by Lean's \texttt{Graph.NodeId}, the JSON graph, and executable traces.  Certification nodes have one successor.  The active ledger is an exact CT4 input; every dormant exit retains the first-failure state and the exact consumer input.}
\label{fig:ct6-active-dormant}
\end{figure}
\FloatBarrier
